\section*{3. Input-output linearization of the mechanics}
Let $q=[\psi,\theta]^\mathsf{T}$ denote yaw and pitch. With perfect rotor-speed tracking, define $s(\omega_d)=[\omega_{m,d}|\omega_{m,d}|,\omega_{t,d}|\omega_{t,d}|]^\mathsf{T}$. The mechanical state-space model is
\begin{equation}
x=\begin{bmatrix}q\\\dot q\end{bmatrix},\qquad
\boxed{\dot x=
\begin{bmatrix}\dot q\\-M_q^{-1}(c_q+g_q)\end{bmatrix}
+\begin{bmatrix}0_{2\times2}\\M_q^{-1}\Gamma(\theta)\end{bmatrix}s(\omega_d)},\qquad
y=\begin{bmatrix}I_2&0_{2\times2}\end{bmatrix}x.
\end{equation}
Using the supplied model, with $\Delta J=J_\psi-j_\psi$,
\begin{equation}
\begin{aligned}
M_q&=\operatorname{diag}(J_\psi\cos^2\theta+j_\psi\sin^2\theta,\ J_\theta),\\
c_q&=\begin{bmatrix}-\Delta J\sin(2\theta)\dot\theta\dot\psi\\
\tfrac12\Delta J\sin(2\theta)\dot\psi^2\end{bmatrix},\qquad
g_q=\begin{bmatrix}0\\G\sin\theta\end{bmatrix},\\
\Gamma(\theta)&=\begin{bmatrix}\cos\theta&0\\0&1\end{bmatrix}
\begin{bmatrix}-k_t&Lk_f\\Lk_f&k_t\end{bmatrix}.
\end{aligned}
\end{equation}
Introduce the virtual body-acceleration input $a\in\mathbb R^2$. The exact inverse is
\begin{equation}
\boxed{s_d=\Gamma(\theta)^{-1}(M_qa+c_q+g_q),\qquad
\omega_d=\operatorname{sign}(s_d)\sqrt{|s_d|},}
\end{equation}
where the square root and sign act componentwise. Substitution gives $\ddot q=a$. The state coordinates remain $z=x=[\psi,\theta,\dot\psi,\dot\theta]^\mathsf{T}$:
\begin{equation}
\dot z=\begin{bmatrix}0_{2\times2}&I_2\\0_{2\times2}&0_{2\times2}\end{bmatrix}z
+\begin{bmatrix}0_{2\times2}\\I_2\end{bmatrix}a,\qquad
y=\begin{bmatrix}I_2&0_{2\times2}\end{bmatrix}z.
\end{equation}
With $s$ as input, each angle has relative degree two; no internal mechanical dynamics remain. Since $\det\Gamma=-\cos\theta[k_t^2+(Lk_f)^2]$, the inverse requires $\cos\theta\ne0$. The matrix $M_q$ is always invertible, since $J_\psi\cos^2\theta+j_\psi\sin^2\theta\ge j_\psi>0$. The simulation restricts $|\theta|\le60^\circ$, inside $|\theta|<90^\circ$. The signed square root is defined at zero but is not differentiable there; a smooth local inverse requires nonzero rotor commands. At steady speed, the voltage limit gives $|\omega_d|\le V_{\max}b_m/a_m=400.8$ rad/s for 24 V. This is a steady-state bound, not a guarantee of transient tracking. Exact cancellation also assumes accurate parameters and measured states. The assumption $\omega\approx\omega_d$ needs speed loops much faster than the position loop.

\section*{4. Linear yaw/pitch position control}
For a desired configuration $q_d$, use measured rates and reference-velocity feedforward:
\begin{equation}
a=K_q(q_d-q)+D_q(\dot q_d-\dot q),\qquad
\ddot e+D_q\dot e+K_qe=\ddot q_d,\quad e=q_d-q.
\end{equation}
Matching each axis polynomial to $(s+\lambda)^2$ gives $K_q=\lambda^2I_2$, $D_q=2\lambda I_2$. For a relative band $\varepsilon$, choose $\lambda=\tau_\varepsilon/T_s$, where $(1+\tau_\varepsilon)e^{-\tau_\varepsilon}=\varepsilon$. For implementation, $\varepsilon=0.02$ and $T_s=3$ s give $\lambda=1.9446$ s$^{-1}$ and
\begin{equation}
\boxed{K_q=\lambda^2I_2=3.7816I_2,\qquad
D_q=2\lambda I_2=3.8893I_2.}
\end{equation}
The supplied controller blocks differentiate the linearly interpolated smooth reference signals to obtain $\dot q_d$; measured angles are not differentiated. Without feedforward, the error equation also contains $D_q\dot q_d$, increasing tracking lag. Sharp-step tests disable reference differentiation and settle in 2.9960 s without overshoot. The laboratory controller instead differentiates the measured position error (see Experimental validation).
\begin{center}
\includegraphics[width=\linewidth]{report_part2_position_steps.png}
\reportcaption{Simulink yaw/pitch step responses with ideal motor tracking. Shading marks the $\pm2\%$ bands; vertical lines mark the 3 s target. Both axes settle in 2.996 s without overshoot.}
\end{center}
